\subsection{Dynamic Sets}

Data structures are used to implement \textbf{dynamic sets}.
Operations are grouped into two categories:

\begin{itemize}
  \item \textbf{Modification}: Insert, Delete, \dots
  \item \textbf{Query}: Search, Minimum, Maximum, \dots
\end{itemize}

% ─────────────────────────────────────────────────────────────────────────────
\subsection{Stacks -- LIFO}

A stack follows the \textbf{LIFO} (\textit{Last-In, First-Out}) principle: last in,
first out, like a stack of plates.

\begin{center}
\begin{tikzpicture}[font=\small, >=Stealth]
  \foreach \v/\i in {3/1, 7/2, 1/3}{
    \node[draw, fill=blue!10, minimum width=1.2cm, minimum height=0.55cm]
          at (0, \i*0.6-0.3) {\v};
  }
  \node[draw, fill=orange!30, minimum width=1.2cm, minimum height=0.55cm]
        at (0, 4*0.6-0.3) {5};
  \draw[->, thick, keyred] (1.2, 4*0.6-0.3) -- (0.65, 4*0.6-0.3)
        node[left=0pt]{}; 
  \node[right=2pt, keyred] at (1.2, 4*0.6-0.3) {\textsc{Push}/\textsc{Pop}};
  \node[left=4pt] at (-0.6, 1) {\textit{bottom}};
  \node[left=4pt] at (-0.6, 4*0.6-0.3) {\textit{top}};
  \node[right=2pt, gray] at (0.65, 0.5) {$S.\mathit{top}$};
\end{tikzpicture}
\end{center}

\begin{itemize}
  \item \textsc{Push}$(S, x)$: pushes $x$ onto the top. $O(1)$.
  \item \textsc{Pop}$(S)$: removes and returns the top element. $O(1)$.
  \item \textsc{Stack-Empty}$(S)$: tests if $S.\mathit{top} = 0$. $O(1)$.
\end{itemize}

% ─────────────────────────────────────────────────────────────────────────────
\subsection{Queues -- FIFO}

A queue follows the \textbf{FIFO} (\textit{First-In, First-Out}) principle: first in,
first out, like a waiting line.

\begin{center}
\begin{tikzpicture}[font=\small, >=Stealth]
  \foreach \v/\i in {3/0, 7/1, 1/2, 9/3}{
    \node[draw, fill=orange!15, minimum width=0.9cm, minimum height=0.65cm]
          at (\i*1.0, 0) {\v};
  }
  \draw[->, thick, keygreen] (-0.9, 0) -- (-0.5, 0)
        node[left=2pt, keygreen]{\textsc{Enqueue}};
  \draw[->, thick, keyred] (3.5+0.5, 0) -- (3.5+0.9, 0)
        node[right=2pt, keyred]{\textsc{Dequeue}};
  \node[below=4pt] at (0, -0.35) {\small$Q.\mathit{head}$};
  \node[below=4pt] at (3, -0.35) {\small$Q.\mathit{tail}$};
\end{tikzpicture}
\end{center}

\begin{itemize}
  \item \textsc{Enqueue}$(Q, x)$: adds $x$ to the tail. $O(1)$.
  \item \textsc{Dequeue}$(Q)$: removes and returns the head element. $O(1)$.
\end{itemize}

% ─────────────────────────────────────────────────────────────────────────────
\subsection{Linked Lists}

Linear order is determined by \textbf{pointers} stored in each object
(and not by array indices).

\begin{center}
\begin{tikzpicture}[font=\small, >=Stealth,
    node/.style={draw, fill=blue!8, minimum width=1.8cm, minimum height=0.65cm,
                  rounded corners=2pt}]
  \node[node] (a) at (0,0) {4 $|$ \texttt{->}};
  \node[node] (b) at (2.4,0) {1 $|$ \texttt{->}};
  \node[node] (c) at (4.8,0) {7 $|$ \texttt{->}};
  \node[node] (d) at (7.2,0) {3 $|$ \texttt{NIL}};
  \draw[->] (a.east) -- (b.west);
  \draw[->] (b.east) -- (c.west);
  \draw[->] (c.east) -- (d.west);
  \node[above=4pt of a, font=\tiny\itshape] {head};
\end{tikzpicture}
\end{center}

\begin{center}
\renewcommand{\arraystretch}{1.25}
\begin{tabular}{lll}
\toprule
\textbf{Operation} & \textbf{Time} & \textbf{Remark} \\
\midrule
\textsc{List-Search}$(L, k)$  & $O(n)$ & Sequential traversal \\
\textsc{List-Insert}$(L, x)$  & $O(1)$ & Insertion at the head \\
\textsc{List-Delete}$(L, x)$  & $O(1)$ & If pointer already known (doubly linked list) \\
\bottomrule
\end{tabular}
\end{center}

\textbf{Sentinels:} a dummy object \texttt{L.nil} represents both ends
of the list simultaneously, simplifying boundary cases (no \texttt{NULL} testing).

\subsubsection{Example: reversing a list in $\Theta(n)$ time, $\Theta(1)$ space}

\begin{algorithm}[H]
\caption{Reverse-List($L$)}
\begin{algorithmic}[1]
\State $prev \leftarrow \textsc{NIL}$
\State $curr \leftarrow L.head$
\While{$curr \ne \textsc{NIL}$}
    \State $next \leftarrow curr.next$
    \State $curr.next \leftarrow prev$   \AlgComment{reverse the link}
    \State $prev \leftarrow curr$
    \State $curr \leftarrow next$
\EndWhile
\State $L.head \leftarrow prev$
\end{algorithmic}
\end{algorithm}

\begin{center}
\begin{tikzpicture}[font=\small, >=Stealth,
    node/.style={draw, fill=blue!8, minimum width=0.8cm, minimum height=0.55cm}]
  % Before
  \node[left] at (-0.3, 0.9) {\small\textit{Before:}};
  \foreach \v/\i in {4/0, 1/1, 7/2, 3/3}{
    \node[node] (\i) at (\i*1.2, 0.9) {\v};
  }
  \draw[->] (0.east)--(1.west); \draw[->] (1.east)--(2.west);
  \draw[->] (2.east)--(3.west);
  % After
  \node[left] at (-0.3, 0) {\small\textit{After:}};
  \foreach \v/\i in {3/0, 7/1, 1/2, 4/3}{
    \node[node] (\i r) at (\i*1.2, 0) {\v};
  }
  \draw[->] (0r.east)--(1r.west); \draw[->] (1r.east)--(2r.west);
  \draw[->] (2r.east)--(3r.west);
\end{tikzpicture}
\end{center}

% ─────────────────────────────────────────────────────────────────────────────
\subsection{Conclusion on Elementary Structures}

\begin{tcolorbox}[colback=lightyellow, colframe=orange!70!black, fonttitle=\bfseries,
                  title={Strengths and Limitations}]
Stacks, queues, and linked lists excel at $O(1)$ insertions/deletions,
but are \textbf{poor for searching} ($O(n)$). For fast searching,
more sophisticated structures are needed (trees, hash tables).
\end{tcolorbox}